\documentclass[11pt]{article}
\usepackage[margin=1.05in]{geometry}
\usepackage{amsmath,amssymb,amsthm,mathtools}
\usepackage{booktabs,array,longtable,microtype}
\usepackage[hidelinks]{hyperref}
\usepackage[nameinlink,capitalise]{cleveref}
\usepackage[numbers,sort&compress]{natbib}
\usepackage{xcolor}
\usepackage{enumitem}
\usepackage{tikz-cd}
\usepackage{listings}

\newtheorem{theorem}{Theorem}[section]
\newtheorem{proposition}[theorem]{Proposition}
\newtheorem{lemma}[theorem]{Lemma}
\newtheorem{corollary}[theorem]{Corollary}
\theoremstyle{definition}
\newtheorem{definition}[theorem]{Definition}
\newtheorem{conjecture}[theorem]{Conjecture}
\theoremstyle{remark}
\newtheorem{remark}[theorem]{Remark}
\newtheorem{warning}[theorem]{Warning}

\DeclareMathOperator{\Hess}{Hess}
\DeclareMathOperator{\per}{per}
\DeclareMathOperator{\rk}{rank}
\DeclareMathOperator{\tr}{tr}
\DeclareMathOperator{\supp}{supp}
\DeclareMathOperator{\Span}{span}
\newcommand{\C}{\mathbb C}
\newcommand{\Q}{\mathbb Q}
\newcommand{\Sym}{\operatorname{Sym}}
\newcommand{\GL}{\operatorname{GL}}
\newcommand{\SL}{\operatorname{SL}}
\newcommand{\Ddet}{D_{\det,n}}
\newcommand{\Dper}{D_{\mathrm{per},n}}
\newcommand{\status}[1]{\textsf{#1}}

\title{Reverse Hessian Obstructions, Catalecticant Blindness,\\and a Program Toward VP versus VNP}
\author{Sasha Lempers\\\small Independent Researcher}
\date{16 July 2026}

\begin{document}
\maketitle

\begin{abstract}
Let $\det_n$ and $\per_n$ denote the determinant and permanent of a generic
$n\times n$ matrix, and put
\[
 D_{\det,n}=\det\Hess(\det_n),\qquad
 D_{\mathrm{per},n}=\det\Hess(\per_n).
\]
We construct a uniform geometric obstruction, based on the Hessian determinant,
that separates the Hessian images of the determinant and the permanent under
orbit closure: for every $n\geq3$,
\[
 D_{\det,n}\notin\overline{\GL_{n^2}(\C)\cdot D_{\mathrm{per},n}}.
\]
The proof combines an exact power identity for the determinant, an explicit
one-parameter restriction for the permanent, and polystability.  In parallel,
we certify that ordinary catalecticants do not see this separation for $n=3$:
the complete rank vectors are
\[
(1,9,45,165,270,270,165,45,9,1)
\quad\text{and}\quad
(1,9,45,165,414,414,165,45,9,1).
\]
Two independently implemented engines reconstruct the forms and ranks, and an
exact gcd certificate proves that $D_{\mathrm{per},3}$ is squarefree.

We then address the relation to the padded permanent used in geometric
complexity theory.  For a homogeneous degree-$d$ form $g$ in $N$ variables we
prove
\[
\det\Hess(\ell^k g)
=-\frac{k(k+d-1)}{d-1}\,
 \ell^{k(N+1)-2}g\,\det\Hess(g)
\]
when the Hessian is taken only in the active variables $(\ell,x)$.  In the
standard GCT ambient space, however, $\ell^{m-n}\per_n$ has unused variables
whenever $m>n$, so its full $m^2\times m^2$ Hessian determinant is identically
zero.  This yields a rigorous barrier: the ordinary full Hessian determinant
alone cannot produce the standard padded-permanent lower bound.  No claim that
VP differs from VNP is made.
\end{abstract}

\section{Introduction}

Valiant's permanent versus determinant problem and the separation of VP from
VNP remain central open problems in algebraic complexity
\cite{Valiant1979,BLMW2011}.  Geometric Complexity Theory (GCT) translates a
strengthened form of the problem into non-containment between orbit closures of
the determinant and a padded permanent
\cite{MulmuleySohoni2001,MulmuleySohoni2008}.  This translation makes
representation theory, invariant theory and algebraic geometry available, but
it also exposes exact barriers.  In particular, broad occurrence-obstruction
programs fail in the relevant asymptotic regime
\cite{IP2017,BIP2019}, while multiplicity obstructions can still be genuinely
stronger in other natural geometric settings \cite{DIP2020}.

This paper starts from a finite, unpadded Hessian-image separation and asks the
question that matters for complexity: can it survive the padding and ambient
embeddings that encode determinantal complexity?  The answer for the ordinary
full Hessian determinant is negative.  The positive finite theorem is real and
fully explicit, but the same construction collapses to zero on every genuinely
padded permanent in the standard ambient space.

The exact contribution is therefore two-sided.
\begin{enumerate}[leftmargin=2em]
\item We prove a uniform orbit-closure separation after applying the
Hessian-determinant covariant in the equal-size unpadded model.
\item We certify a complete ordinary-catalecticant no-go in size three and
record finite negative evidence for three shifted-partial-derivative ranks.
\item We prove an active-variable padding factorisation and a full-ambient
padding-collapse theorem.
\item We isolate what an anti-padding covariant would have to accomplish before
a Hessian program could have consequences for VP versus VNP.
\end{enumerate}

The status labels used throughout are:
\status{PROVED} for a complete paper proof;
\status{CERTIFIED-COMP} for exact reproducible computation;
\status{LITERATURE} for an external theorem used with a precise citation;
\status{CONJECTURAL}, \status{FAILED}, and \status{UNKNOWN} otherwise.
Finite checks never replace the general proofs.

\section{Background on VP, VNP, and GCT}

\subsection{Permanent, determinant and padding}

For a generic $n\times n$ matrix $X=(x_{ij})$, write
\[
\det_n(X)=\sum_{\sigma\in S_n}\operatorname{sgn}(\sigma)
 \prod_i x_{i,\sigma(i)},\qquad
\per_n(X)=\sum_{\sigma\in S_n}\prod_i x_{i,\sigma(i)}.
\]
A determinantal representation of a polynomial $f$ of degree $n$ is an affine
linear matrix $A(x)$ of size $m$ such that $f=\det A(x)$.  Determinantal
complexity is the least such $m$; border determinantal complexity allows a
limit of such representations.  The GCT homogenisation uses the degree-$m$
form
\[
\ell^{m-n}\per_n\in\Sym^m(\C^{m^2})^*,
\]
where the $n^2$ permanent variables and a padding variable $\ell$ are embedded
among $m^2$ ambient variables.  The corresponding determinant form is
$\det_m\in\Sym^m(\C^{m^2})^*$.

The equal-size comparison between $\det_n$ and $\per_n$ in $n^2$ variables is
not the standard padded formulation.  Both forms are known to be polystable,
and their stabilizers are classical; see \cite{BI2017}.  Kumar proves
nonnormality phenomena for determinant and padded-permanent orbit closures
\cite{Kumar2013}, and the boundary of the $3\times3$ determinant orbit is
explicitly described in \cite{HuttenhainLairez2016}.  None of these facts may
be replaced by an unsupported normality or openness assumption.

\subsection{Logical map to VP versus VNP}

The relevant implications must be separated by computational model.  The
permanent family is VNP-complete under p-projections, so
\[
 \per\notin\mathrm{VP}\quad\Longleftrightarrow\quad
 \mathrm{VP}\ne\mathrm{VNP}.
\]
Polynomial determinantal complexity instead characterises the weakly-skew
class $\mathrm{VP}_{\mathrm{ws}}$ (equivalently, polynomial-size algebraic
branching programs up to standard polynomial simulations).  Consequently, a
super-polynomial lower bound for determinantal complexity of the permanent
would prove
$\mathrm{VP}_{\mathrm{ws}}\ne\mathrm{VNP}$, a major lower bound but not, by
itself, the full separation $\mathrm{VP}\ne\mathrm{VNP}$; see
\cite{BlaserIkenmeyer2025}.

For the border model, the standard GCT implication has the following direction:
\[
 \ell^{m-n}\per_n\in\overline{\GL_{m^2}\cdot\det_m}
 \quad\Longrightarrow\quad
 C(\ell^{m-n}\per_n)
 \in\overline{\GL\cdot C(\det_m)}                         \tag{2.1}
\]
for a regular equivariant covariant $C$ with the appropriate character twist.
Thus the complexity-relevant target obstruction would be
\[
 C(\ell^{m-n}\per_n)\notin\overline{\GL\cdot C(\det_m)}. \tag{2.2}
\]
The theorem proved in this paper has the opposite orientation,
$D_{\det,n}\notin\overline{\GL\cdot D_{\mathrm{per},n}}$, and is therefore
not a GCT lower-bound certificate even before padding is introduced.  After
padding, the full Hessian determinant has the stronger defect
$C(\ell^{m-n}\per_n)=0$, while $0$ lies in the closure of every nonzero
positive-degree homogeneous orbit.  The status map is therefore:
\begin{center}
\begin{tabular}{p{0.54\textwidth}p{0.28\textwidth}}
\toprule
Statement & Status\\\midrule
super-polynomial unrestricted circuit lower bound for $\per$ $\Rightarrow$ $\mathrm{VP}\ne\mathrm{VNP}$ & known\\
super-polynomial determinantal complexity of $\per$ $\Rightarrow$ $\mathrm{VP}_{\mathrm{ws}}\ne\mathrm{VNP}$ & known\\
padded orbit-closure non-containment for all polynomial $m(n)$ $\Rightarrow$ border determinantal lower bound & known in the precise GCT model\\
reverse unpadded Hessian-image non-containment $\Rightarrow$ padded permanent non-containment & false direction / missing\\
ordinary full Hessian determinant $\Rightarrow$ padded lower bound & impossible by \cref{thm:collapse}\\
anti-padding transverse covariant $\Rightarrow$ standard GCT lower bound & unknown\\
\bottomrule
\end{tabular}
\end{center}

\section{The Reverse Hessian Obstruction}

Let $W$ be an $M$-dimensional complex vector space and
$f\in\Sym^dW^*$.  We use
\[
(g\cdot f)(x)=f(g^{-1}x),\qquad g\in\GL(W),
\]
and write $\mathcal D(f)=\det\Hess(f)$ in fixed linear coordinates.

\begin{lemma}[Hessian covariance]\label{lem:covariance}
For every $g\in\GL(W)$,
\[
\mathcal D(g\cdot f)=\det(g)^{-2}\,g\cdot\mathcal D(f).
\]
In particular, $\mathcal D$ is $\SL(W)$-equivariant.
\end{lemma}
\begin{proof}
The chain rule gives
$\Hess(g\cdot f)(x)=g^{-T}\Hess(f)(g^{-1}x)g^{-1}$.
Taking determinants proves the claim.
\end{proof}

\subsection{Exact determinant identity}

\begin{theorem}\label{thm:det-hessian}
For every $n\ge2$,
\[
\det\Hess(\det_n)=
(n-1)(-1)^{(n-1)(n+2)/2}\det_n^{\,n(n-2)}.
\]
In particular, $\det\Hess(\det_3)=-2\det_3^3$.
\end{theorem}
\begin{proof}
For $X\in\GL_n$ and $U,V\in M_n$, Jacobi's formula gives
\[
d^2(\det)_X(U,V)=\det(X)\bigl(
\tr(X^{-1}U)\tr(X^{-1}V)-\tr(X^{-1}UX^{-1}V)\bigr).
\]
At $I_n$ the Hessian bilinear form is
$B(U,V)=\tr(U)\tr(V)-\tr(UV)$.  At general $X$, the Hessian is
$\det(X)$ times the pullback of $B$ by left multiplication
$L_{X^{-1}}$.  Since $\det(L_{X^{-1}})=\det(X)^{-n}$ and the Hessian
has size $n^2$,
\[
\det\Hess(\det_n)(X)=\det(B)\det(X)^{n^2-2n}.
\]
The identity extends polynomially from $\GL_n$ to $M_n$.

On the diagonal subspace, $B$ has matrix $J_n-I_n$ and determinant
$(n-1)(-1)^{n-1}$.  For every unordered pair $i<j$, the span of
$E_{ij},E_{ji}$ contributes the block
$\left(\begin{smallmatrix}0&-1\\-1&0\end{smallmatrix}\right)$,
with determinant $-1$.  Multiplying the $n(n-1)/2$ off-diagonal blocks gives
\[
\det(B)=(n-1)(-1)^{n-1+n(n-1)/2}
=(n-1)(-1)^{(n-1)(n+2)/2}.
\]
\end{proof}
The power shape is compatible with the general theory of relative invariants
and Hessian determinants \cite{Fox2016}; the proof above fixes the exact scalar.

\subsection{A permanent line restriction}

Let $X(t)=J_n+(t-1)E_{11}$, $A=J_n-I_n$, and
$C=0\oplus(J_{n-1}-I_{n-1})$.

\begin{theorem}[Exact line restriction]\label{thm:per-line}
For every $n\ge3$,
\[
\det\Hess(\per_n)(X(t))=
\frac{((n-2)!)^{n^2}(n-1)^{2n}}{(n-2)^{(n-2)^2}}
(t+n-3)^{(n-2)^2}.
\]
\end{theorem}
\begin{proof}
Index the Hessian by pairs $(i,j)$.  Multilinearity gives zero whenever two
differentiated variables share a row or column.  Otherwise the entry is the
permanent of the complementary $(n-2)\times(n-2)$ minor.  At $X(t)$ this is
$(n-2)!$, except when neither row $1$ nor column $1$ is removed, where it is
$(n-3)!(t+n-3)$.  Hence
\[
\Hess(\per_n)(X(t))=(n-2)!A\otimes A+(t-1)(n-3)!C\otimes C.
\]
Put $B=A^{-1}C$.  Since
$A^{-1}=-I_n+\frac1{n-1}J_n$, the map $B$ is the identity on
\[
U_0=\{(0,u_2,\dots,u_n):u_2+\cdots+u_n=0\},
\]
of dimension $n-2$, and is nilpotent on
$\langle e_1,e_2+\cdots+e_n\rangle$.  Its eigenvalues are therefore $1$ with
multiplicity $n-2$ and $0$ with algebraic multiplicity $2$.  Factoring the
Kronecker expression and setting $\alpha=(t-1)/(n-2)$ yields
\begin{align*}
\det\Hess(\per_n)(X(t))
&=((n-2)!)^{n^2}\det(A\otimes A)\det(I+\alpha B\otimes B)\\
&=((n-2)!)^{n^2}(n-1)^{2n}(1+\alpha)^{(n-2)^2},
\end{align*}
which is the stated formula.
\end{proof}

\begin{corollary}\label{cor:not-power}
For $n\ge3$, $D_{\mathrm{per},n}$ is not a nonzero scalar multiple of an
$n(n-2)$-th power of a degree-$n$ form.
\end{corollary}
\begin{proof}
An $r$-th power has every nonzero line-restriction root with multiplicity
divisible by $r=n(n-2)$.  The restriction in \cref{thm:per-line} has a root of
multiplicity $(n-2)^2$, and
$0<(n-2)^2<n(n-2)$.
\end{proof}

\subsection{Polystability and closure}

We use the support-cone criterion of
\cite[Proposition~2.8]{BI2017}.

\begin{proposition}\label{prop:poly}
For every $n\ge3$, $D_{\det,n}$ and $D_{\mathrm{per},n}$ are
$\SL(M_n)$-polystable.
\end{proposition}
\begin{proof}
Let $R$ be the row-column diagonal torus
$E_{ij}\mapsto a_ib_jE_{ij}$ with
$\prod_i a_i=\prod_j b_j=1$.  It lies in $\SL(M_n)$ and stabilizes both
$\det_n$ and $\per_n$, hence both Hessian forms by \cref{lem:covariance}.
Its weights on the one-dimensional spaces $\C E_{ij}$ are pairwise distinct,
so its centralizer is diagonal.

The group $S_n\times S_n$ acts transitively on the $n^2$ variables and
preserves each Hessian support.  Both supports are nonempty by
\cref{thm:det-hessian,thm:per-line}.  Each Hessian form has degree
$q=n^2(n-2)$.  Averaging any support exponent under the transitive permutation
group gives $(n-2)(1,\ldots,1)$; therefore the support cone contains the
all-ones vector.  The cited criterion gives polystability.
\end{proof}

\begin{lemma}[Homogeneous GL closure lemma]\label{lem:glsl}
Let $U$ be a homogeneous polynomial representation of $\GL_M$ whose center
generates nonzero scalar multiplication.  If nonzero $u,v\in U$ are
$\SL_M$-polystable and $u\in\overline{\GL_M\cdot v}$, then
$u\in\GL_M\cdot v$.
\end{lemma}
\begin{proof}
Write $\GL_M\cdot v=\C^*(\SL_M\cdot v)$ and choose
$c_j\in\C^*$, $g_j\in\SL_M$ with $c_jg_jv\to u$.  Since $u$ is nonzero and
polystable, an $\SL_M$-invariant homogeneous polynomial $P$ of positive degree
satisfies $P(u)\ne0$.  Then
$c_j^{\deg P}P(v)\to P(u)$, so $P(v)\ne0$ and, after taking a subsequence,
$c_j\to c\ne0$.  Thus $g_jv\to c^{-1}u$.  The $\SL_M$ orbit of $v$ is closed,
so $c^{-1}u\in\SL_M\cdot v$.
\end{proof}

\begin{theorem}[Uniform reverse-Hessian separation]\label{thm:reverse}
For every $n\ge3$,
\[
D_{\det,n}\notin\overline{\GL(M_n)\cdot D_{\mathrm{per},n}}.
\]
\end{theorem}
\begin{proof}
By \cref{thm:det-hessian}, $D_{\det,n}$ is an $n(n-2)$-th power up to a
nonzero scalar.  By \cref{cor:not-power}, $D_{\mathrm{per},n}$ is not.
Thus they are not in the same $\GL(M_n)$ orbit.  If the displayed containment
held, \cref{prop:poly,lem:glsl} would force orbit equivalence, a contradiction.
\end{proof}

\begin{warning}[Direction mismatch]
\Cref{thm:reverse} concerns Hessian images in an equal-size unpadded model and
has the reverse orientation from the complexity-relevant obstruction in
\textup{(2.2)}.  It neither proves a non-containment for the original padded
permanent nor gives a circuit, determinantal-complexity, or border-
determinantal-complexity lower bound.
\end{warning}

\section{Reproducible Certification}

The new certification chain does not trust historical JSON outputs.  It
reconstructs all relevant polynomials from monomial data and separates
mathematical certificates from machine-dependent runtime logs.

\subsection{Two independent catalecticant engines}

The primary engine builds $\det_3$, $\per_3$, their $9\times9$ Hessian
determinants, and ordinary apolar catalecticant matrices.  It computes ranks for
orders $0$ through $4$ exactly over $\Q$ using sparse domain matrices; orders
$5$ through $9$ follow from the proved transpose/factorial duality.

The clean-room engine imports no code from the primary engine.  It represents
polynomials as exponent dictionaries, differentiates them directly,
reconstructs both Hessian determinants, builds all ten matrices, and performs
its own sparse Gaussian elimination modulo
$1009$, $32003$, and $65521$.  The orchestration script rejects the run unless
polynomial hashes, direct matrix hashes and every modular rank agree with the
primary exact rank.

\begin{table}[ht]
\centering
\small
\begin{tabular}{c@{\qquad}r@{\qquad}r}
\toprule
$a$ & $\rk C_a(D_{\det,3})$ & $\rk C_a(D_{\mathrm{per},3})$\\
\midrule
0&1&1\\1&9&9\\2&45&45\\3&165&165\\4&270&414\\
5&270&414\\6&165&165\\7&45&45\\8&9&9\\9&1&1\\
\bottomrule
\end{tabular}
\caption{Exact ordinary catalecticant rank vectors.}
\label{tab:ranks}
\end{table}

\subsection{Squarefreeness}

\begin{lemma}[Derivative gcd criterion]\label{lem:squarefree}
Let $K$ be a characteristic-zero field and $0\ne F\in K[x_1,\ldots,x_N]$.
If $\gcd(\partial_1F,\ldots,\partial_NF)=1$, then $F$ is squarefree.
\end{lemma}
\begin{proof}
If an irreducible polynomial $h$ occurred in $F$ with multiplicity at least
$2$, writing $F=h^2q$ would show that $h$ divides every first derivative of
$F$, contradicting the gcd hypothesis.
\end{proof}

A separate deterministic script reconstructs $D_{\mathrm{per},3}$ over
$\mathbb Z$, computes all nine first partial derivatives, and obtains gcd $1$
over $\Q$.  By \cref{lem:squarefree}, this certifies squarefreeness over $\Q$.
Squarefreeness is preserved by the field extension $\Q\subset\C$, hence also
holds over $\C$.  The integer content gcd of the partial derivatives is $2$;
the modular data are sanity checks, not the proof.

\subsection{Bit-for-bit reproducibility}

Certificates contain no absolute paths, elapsed times, hostnames, or timestamps.
Running
\begin{lstlisting}[language=bash]
python run_all.py
python scripts/check_reproducibility.py
\end{lstlisting}
rebuilds the certificates in a fresh temporary directory and compares all eight
certificate files byte for byte.  \path{certificates/SHA256SUMS.txt} hashes the
mathematical outputs; runtime metadata is isolated under \path{runtime_logs/}.

\section{Catalecticant Blindness}

For a degree-nine form $F\in\Sym^9V^*$, the ordinary apolar catalecticant is
\[
C_a(F):\Sym^aV\longrightarrow\Sym^{9-a}V^*,\qquad
\theta\longmapsto\theta\cdot F.
\]
It is equivariant, and its rank is upper semicontinuous.

\begin{lemma}[Catalecticant duality]\label{lem:catduality}
For every degree-$d$ form $F$ and every $0\le a\le d$,
\[
 \rk C_a(F)=\rk C_{d-a}(F).
\]
\end{lemma}
\begin{proof}
Use divided-power monomial bases in $\Sym^aV$ and $\Sym^{d-a}V$.
The coefficient of a monomial indexed by $\beta$ in
$C_a(F)(\partial^\alpha)$ is the coefficient of the degree-$d$ monomial
indexed by $\alpha+\beta$ in $F$.  The same coefficient is the
$(\alpha,\beta)$ entry of the transpose matrix for $C_{d-a}(F)$.  In ordinary
monomial bases the two matrices differ only by invertible diagonal factorial
scalings.  Their ranks are therefore equal.
\end{proof}

\begin{theorem}[Complete ordinary-catalecticant no-go]\label{thm:catblind}
Every ordinary catalecticant rank condition that contains
$\overline{\GL(V)\cdot D_{\mathrm{per},3}}$ also contains
$D_{\det,3}$.
\end{theorem}
\begin{proof}
For each $a$, the closed rank locus containing the permanent orbit closure has
threshold at least $\rk C_a(D_{\mathrm{per},3})$.  The certified table gives
\[
\rk C_a(D_{\det,3})\le\rk C_a(D_{\mathrm{per},3})
\quad(0\le a\le9).
\]
Thus $D_{\det,3}$ lies in every such locus and in every intersection of them.
\end{proof}

This theorem concerns rank conditions only.  It does not rule out non-rank
equations in catalecticant entries, Young or Koszul flattenings, apolar
syzygies, singular-locus invariants, or multiplicity obstructions.  Such
constructions are central in the geometry of secant and complexity varieties
\cite{LandsbergOttaviani2013,Guan2017}.

There is also a distinct asymptotic limitation in the literature:
Efremenko--Landsberg--Schenck--Weyman prove that shifted partial derivatives
alone cannot separate the standard padded permanent from the determinant orbit
closure in their stated parameter range \cite{ELSW2016}.  The finite tests below
concern the two Hessian-image forms and neither reprove nor strengthen that
result.

Three exact shifted-partial tests were also run.  Their ranks are
\[
\begin{array}{c|cc}
(k,\ell)&D_{\det,3}&D_{\mathrm{per},3}\\\hline
(1,1)&65&77\\
(2,1)&270&369\\
(1,2)&270&369
\end{array}
\]
and again have the wrong inequality for the desired reverse separation.  These
are finite \status{CERTIFIED-COMP} failures, not a general theorem about shifted
partial derivatives.

\section{Relation to Padded Permanent Orbit Closures}

\subsection{What covariance does and does not imply}

Because $\mathcal D$ is a polynomial covariant up to a determinant character,
an original orbit containment forces a target containment after the appropriate
twist.  A target non-containment can therefore obstruct an original
containment when both target values are meaningful.  No converse holds: target
containment, equality, or vanishing need not reflect the original geometry.

The standard padded permanent exhibits the worst possible loss: its full
Hessian determinant is zero.

\subsection{Exact active-variable padding formula}

\begin{theorem}[One-padding-variable Hessian factorisation]\label{thm:paddingfactor}
Let $g\in\C[x_1,\ldots,x_N]$ be homogeneous of degree $d>1$, let $k\ge1$, and
put $F(\ell,x)=\ell^kg(x)$.  Then
\[
\det\Hess_{(\ell,x)}(F)
=-\frac{k(k+d-1)}{d-1}\,
\ell^{k(N+1)-2}g(x)\det\Hess_x(g).
\]
The identity holds as a polynomial even when $\Hess_x(g)$ is generically
singular.
\end{theorem}
\begin{proof}
Put $H=\Hess_x(g)$ and $u=\nabla g$.  The full Hessian has block form
\[
\begin{pmatrix}
k(k-1)\ell^{k-2}g & k\ell^{k-1}u^T\\
k\ell^{k-1}u & \ell^kH
\end{pmatrix}.
\]
For a scalar $a$, a column $b$ and a square matrix $D$, the polynomial block
identity
\[
\det\begin{pmatrix}a&b^T\\b&D\end{pmatrix}
=a\det D-b^T\operatorname{adj}(D)b
\]
holds without any invertibility hypothesis.  Factoring the powers of $\ell$
therefore gives
\[
\det\Hess(F)=\ell^{k(N+1)-2}
\left(k(k-1)g\det H-k^2u^T\operatorname{adj}(H)u\right).
\]
Differentiating Euler's identity $x^Tu=dg$ gives $Hx=(d-1)u$.
Multiplication by $\operatorname{adj}(H)$ yields
\[
\det(H)x=(d-1)\operatorname{adj}(H)u.
\]
Taking the scalar product with $u$ and using $u^Tx=dg$ gives the polynomial
identity
\[
u^T\operatorname{adj}(H)u=\frac{d}{d-1}g\det H.
\]
Substitution now gives
\[
 k(k-1)-\frac{k^2d}{d-1}
 =-\frac{k(k+d-1)}{d-1},
\]
which proves the formula everywhere, including when $H$ is generically
singular.
\end{proof}

This formula recovers a canonical trace of $g\det\Hess(g)$ only when the padding
line and the active quotient are already marked.  It is not, by itself, a
regular $\GL$-covariant on unmarked forms.

\subsection{Full-ambient collapse}

\begin{lemma}[Unused-variable collapse]\label{lem:unused}
If a polynomial in $M$ ambient variables is independent of one variable, its
full $M\times M$ Hessian determinant is zero.
\end{lemma}
\begin{proof}
The row and column indexed by that variable are zero.
\end{proof}

\begin{theorem}[Padding-collapse barrier]\label{thm:collapse}
Let $m>n$ and regard
$F_{n,m}=\ell^{m-n}\per_n$ as a degree-$m$ form in $m^2$ variables in the
standard padded GCT model.  Then
\[
\det\Hess_{\C^{m^2}}(F_{n,m})=0.
\]
Consequently, the ordinary full Hessian-determinant covariant alone cannot
separate the padded permanent orbit closure from the determinant orbit closure
or imply an asymptotic determinantal-complexity lower bound.
\end{theorem}
\begin{proof}
The form $F_{n,m}$ depends on at most $n^2+1$ variables: the $n^2$ permanent
variables and $\ell$.  Since
\[
m^2-(n^2+1)=(m-n)(m+n)-1>0
\]
for $m>n$, there is an unused ambient variable.  The determinant vanishes by
\cref{lem:unused}.

The image of the padded permanent is therefore the zero form.  For any nonzero
positive-degree form $H$, the center of $\GL_{m^2}$ scales $H$, and zero belongs
to $\overline{\GL_{m^2}\cdot H}$.  Thus the zero target cannot exclude the
padded permanent from any determinant target orbit closure.  A lower bound
would need a different covariant or a canonical active-space refinement.
\end{proof}

\section{Beyond Ordinary Catalecticants}

The barrier does not rule out differential covariants that retain all Hessian
minors, the polar map, Fitting ideals of the derivative module, or a marked
transverse Hessian.  It only rules out the scalar determinant of the full
ambient Hessian.

A viable anti-padding construction must satisfy all of the following:
\begin{enumerate}[leftmargin=2em]
\item be regular and equivariant on a space stable under the relevant
$\GL_{m^2}$ action, or explicitly work in a restricted marked model;
\item survive unused ambient variables without making a noncanonical choice;
\item retain a permanent-side invariant that scales uniformly with $n,m$;
\item constrain determinant images by a theorem that is stable under orbit
closure;
\item evade the known occurrence-obstruction limitations
\cite{IP2017,BIP2019}; and
\item be connected by a proved reduction to a standard complexity model.
\end{enumerate}

A first concrete object is the incidence variety of pairs $(F,U)$ with
$F\in\Sym^mU^*$ and $U\subset\C^{m^2}$ an active subspace.  On a fixed stratum,
the maximal nonzero Hessian minors are polynomial, and
\cref{thm:paddingfactor} identifies their marked padding behaviour.  The
unresolved issue is descent: the minimal active subspace can jump in a
 degeneration, and choosing it need not be regular.

\section{A Candidate Asymptotic Obstruction / Barrier Theorem}

The publishable formal result produced by the upgrade is the negative theorem
\cref{thm:collapse}.  It is more useful than an unsupported positive claim
because it removes the ordinary full Hessian determinant from the list of
possible direct GCT covariants for the padded model.

\begin{conjecture}[Canonical active-space extraction]\label{conj:active}
There exists a finite regular equivariant incidence construction whose value on
padded-permanent orbits determines the active subspace and the transverse
Hessian data of the induced active form.
\end{conjecture}
This conjecture is not equivalent to VP versus VNP.  Even if proved, one would
still need an asymptotic determinant-versus-permanent separation in its target
and a reduction back to the standard unmarked model.

\begin{conjecture}[Marked transverse-Hessian multiplicity gap]
In a suitable parabolic model retaining the padding line, an explicit family of
highest weights has a multiplicity gap between transverse-Hessian images of the
padded permanent and determinant families.
\end{conjecture}
No partition is presently certified.  The conjecture is motivated only by the
fact that multiplicity obstructions can be stronger than occurrence
obstructions in other natural varieties \cite{DIP2020}.  Its plausibility is
low until a size-$3$ decomposition is computed exactly.

\section{Consequences and Limitations}

\subsection{Exact consequences}

\begin{itemize}[leftmargin=2em]
\item The reverse Hessian images are separated uniformly in the unpadded
same-size model; this is explicitly the opposite orientation from the
complexity-relevant padded obstruction.
\item All ordinary catalecticant rank tests are blind to that separation for
$n=3$.
\item The standard padded permanent is annihilated by the full ordinary Hessian
determinant for every genuine padding $m>n$.
\item A marked active-variable Hessian retains the factor
$g\det\Hess(g)$ with an explicit scalar and padding exponent.
\end{itemize}

\subsection{What is not proved}

There is no proof of VP versus VNP, no separation of
$\mathrm{VP}_{\mathrm{ws}}$ from VNP, no unrestricted arithmetic-circuit lower
bound, no super-polynomial determinantal- or border-determinantal-complexity
lower bound, no padded orbit
non-containment obtained from the full Hessian determinant, no Young/Koszul
separation, and no coordinate-ring multiplicity gap.  The historical degree-six
C0 sweep is archived but is not promoted into the claims of this paper.

Recent unconditional lower bounds for symmetric arithmetic circuits show that
natural restricted models can be separated: the permanent has exponential
lower bounds in the symmetric model of \cite{DawarWilsenach2025}, and the
homomorphism-polynomial characterisations of
\cite{DPS2026,DwivediPagoSeppelt2026} develop a broader symmetric algebraic
complexity theory.  These results do not imply unrestricted VP versus VNP, and
the present Hessian barrier should be compared to them only after a precise
restricted model is defined.  Algebraic-natural-proofs barriers provide a
second warning against promoting an efficiently checkable finite property to a
general lower-bound method without testing its largeness and usefulness
\cite{GKSS2017}.

\appendix
\section{Exact Scripts and Certificates}

The active files are:
\begin{longtable}{p{0.36\textwidth}p{0.56\textwidth}}
\toprule
File & Role\\\midrule
\path{scripts/certify_catalecticant_ranks.py} & exact rational primary engine\\
\path{scripts/independent_catalecticant_check.py} & independent clean-room modular engine\\
\path{scripts/certify_squarefree_Dper3.py} & exact gcd certificate\\
\path{scripts/verify_uniform_formulas.py} & finite implementation checks\\
\path{scripts/transverse_hessian_experiments.py} & exact padding-factor checks\\
\path{scripts/exploratory_flattenings.py} & finite shifted-partial ranks\\
\path{scripts/check_reproducibility.py} & isolated byte-for-byte rebuild\\
\path{certificates/SHA256SUMS.txt} & deterministic output manifest\\
\bottomrule
\end{longtable}

The exact certificate hashes are recorded in
\path{research_notes/RESULT_REPRODUCTION_TABLE.md}.  The complete 202-file source
archive inventory is in \path{research_notes/AUDIT_ARCHIVE.md}.

\section{Literature Audit}

The literature review was performed before assigning any novelty status.
The distinction between VP, weakly-skew VP, determinantal complexity and the
padded border model follows \cite{BlaserIkenmeyer2025}.  GCT
I/II establish the orbit-closure program \cite{MulmuleySohoni2001,MulmuleySohoni2008};
\cite{IP2017,BIP2019} rule out broad occurrence-obstruction strategies,
while \cite{DIP2020} shows that multiplicity obstructions can be strictly
stronger in natural Chow/secant settings; \cite{BI2017} supplies the polystability/support criterion and fundamental
invariant context; \cite{Kumar2013} and \cite{HuttenhainLairez2016} constrain
claims about determinant and padded-permanent orbit closures; and
\cite{KadishLandsberg2014} studies padded polynomial varieties directly.
The shifted-partial limitation \cite{ELSW2016}, the algebraic-natural-proofs
program \cite{GKSS2017}, and the recent symmetric-circuit results
\cite{DawarWilsenach2025,DPS2026,DwivediPagoSeppelt2026} delimit the relevant
barrier and restricted-model landscape.

No searched source was found to state the exact combination of
\cref{thm:per-line,thm:reverse,thm:catblind,thm:collapse}.  Search failure is
not proof of novelty.  The absolute novelty status is therefore
\status{UNKNOWN} pending specialist review and peer review.  Full titles, years,
DOIs and URLs are listed in \path{research_notes/literature_review.md}.

\begin{thebibliography}{99}
\bibitem{Valiant1979} L.~G. Valiant, \emph{Completeness Classes in Algebra}, STOC 1979, 249--261. DOI: \href{https://doi.org/10.1145/800135.804419}{10.1145/800135.804419}.
\bibitem{MulmuleySohoni2001} K.~D. Mulmuley and M. Sohoni, \emph{Geometric Complexity Theory I: An Approach to the P vs. NP and Related Problems}, SIAM J. Comput. 31 (2001), 496--526. DOI: \href{https://doi.org/10.1137/S009753970038715X}{10.1137/S009753970038715X}.
\bibitem{MulmuleySohoni2008} K.~D. Mulmuley and M. Sohoni, \emph{Geometric Complexity Theory II: Towards Explicit Obstructions for Embeddings among Class Varieties}, SIAM J. Comput. 38 (2008), 1175--1206. DOI: \href{https://doi.org/10.1137/080718115}{10.1137/080718115}.
\bibitem{BLMW2011} P. B{\"u}rgisser, J.~M. Landsberg, L. Manivel and J. Weyman, \emph{An Overview of Mathematical Issues Arising in the Geometric Complexity Theory Approach to VP versus VNP}, SIAM J. Comput. 40 (2011), 1179--1209. DOI: \href{https://doi.org/10.1137/090765328}{10.1137/090765328}.
\bibitem{BlaserIkenmeyer2025} M. Bl{\"a}ser and C. Ikenmeyer, \emph{Introduction to Geometric Complexity Theory}, Theory of Computing Graduate Surveys 10 (2025), 1--166. \url{https://theoryofcomputing.org/articles/gs010/}.
\bibitem{GesmundoIkenmeyerPanova2017} F. Gesmundo, C. Ikenmeyer and G. Panova, \emph{Geometric Complexity Theory and Matrix Powering}, Differential Geom. Appl. 55 (2017), 106--127. DOI: \href{https://doi.org/10.1016/j.difgeo.2017.07.001}{10.1016/j.difgeo.2017.07.001}.
\bibitem{IP2017} C. Ikenmeyer and G. Panova, \emph{Rectangular Kronecker Coefficients and Plethysms in Geometric Complexity Theory}, Adv. Math. 319 (2017), 40--66. DOI: \href{https://doi.org/10.1016/j.aim.2017.08.024}{10.1016/j.aim.2017.08.024}.
\bibitem{BIP2019} P. B{\"u}rgisser, C. Ikenmeyer and G. Panova, \emph{No Occurrence Obstructions in Geometric Complexity Theory}, J. Amer. Math. Soc. 32 (2019), 163--193. DOI: \href{https://doi.org/10.1090/jams/908}{10.1090/jams/908}.
\bibitem{DIP2020} J. D{\"o}rfler, C. Ikenmeyer and G. Panova, \emph{On Geometric Complexity Theory: Multiplicity Obstructions Are Stronger than Occurrence Obstructions}, SIAM J. Appl. Algebra Geom. 4(2) (2020), 354--376. DOI: \href{https://doi.org/10.1137/19M1287638}{10.1137/19M1287638}.
\bibitem{BI2017} P. B{\"u}rgisser and C. Ikenmeyer, \emph{Fundamental Invariants of Orbit Closures}, J. Algebra 477 (2017), 390--434. DOI: \href{https://doi.org/10.1016/j.jalgebra.2016.12.035}{10.1016/j.jalgebra.2016.12.035}.
\bibitem{Kumar2013} S. Kumar, \emph{Geometry of Orbits of Permanents and Determinants}, Comment. Math. Helv. 88 (2013), 759--788. DOI: \href{https://doi.org/10.4171/CMH/302}{10.4171/CMH/302}.
\bibitem{HuttenhainLairez2016} J. H{\"u}ttenhain and P. Lairez, \emph{The Boundary of the Orbit of the 3-by-3 Determinant Polynomial}, C. R. Math. 354 (2016), 931--935. DOI: \href{https://doi.org/10.1016/j.crma.2016.07.002}{10.1016/j.crma.2016.07.002}.
\bibitem{KadishLandsberg2014} H. Kadish and J.~M. Landsberg, \emph{Padded Polynomials, Their Cousins, and Geometric Complexity Theory}, Comm. Algebra 42 (2014), 2171--2180. DOI: \href{https://doi.org/10.1080/00927872.2012.758268}{10.1080/00927872.2012.758268}.
\bibitem{Fox2016} D.~J.~F. Fox, \emph{Functions Dividing Their Hessian Determinants and Affine Spheres}, Asian J. Math. 20 (2016), 503--530. DOI: \href{https://doi.org/10.4310/AJM.2016.v20.n3.a5}{10.4310/AJM.2016.v20.n3.a5}.
\bibitem{LandsbergOttaviani2013} J.~M. Landsberg and G. Ottaviani, \emph{Equations for Secant Varieties of Veronese and Other Varieties}, Ann. Mat. Pura Appl. 192 (2013), 569--606. DOI: \href{https://doi.org/10.1007/s10231-011-0238-6}{10.1007/s10231-011-0238-6}.
\bibitem{Guan2017} Y.~Q. Guan, \emph{Flattenings and Koszul Young Flattenings Arising in Complexity Theory}, Comm. Algebra 45 (2017), 4002--4017. DOI: \href{https://doi.org/10.1080/00927872.2016.1253706}{10.1080/00927872.2016.1253706}.
\bibitem{ELSW2016} K. Efremenko, J. M. Landsberg, H. Schenck and J. Weyman, \emph{The Method of Shifted Partial Derivatives Cannot Separate the Permanent from the Determinant}, arXiv:1609.02103 (2016), DOI: \href{https://doi.org/10.48550/arXiv.1609.02103}{10.48550/arXiv.1609.02103}.
\bibitem{GKSS2017} J.~A. Grochow, M. Kumar, M. Saks and S. Saraf, \emph{Towards an Algebraic Natural Proofs Barrier via Polynomial Identity Testing}, arXiv:1701.01717 (2017), \url{https://arxiv.org/abs/1701.01717}.
\bibitem{DawarWilsenach2025} A. Dawar and G. Wilsenach, \emph{Symmetric Arithmetic Circuits}, Theory of Computing 21(14) (2025). DOI: \href{https://doi.org/10.4086/toc.2025.v021a014}{10.4086/toc.2025.v021a014}.
\bibitem{DPS2026} A. Dawar, B. Pago and T. Seppelt, \emph{Symmetric Algebraic Circuits and Homomorphism Polynomials}, ITCS 2026, LIPIcs 362, pp. 46:1--46:15. DOI: \href{https://doi.org/10.4230/LIPIcs.ITCS.2026.46}{10.4230/LIPIcs.ITCS.2026.46}.
\bibitem{DwivediPagoSeppelt2026} P. Dwivedi, B. Pago and T. Seppelt, \emph{Lower Bounds in Algebraic Complexity via Symmetry and Homomorphism Polynomials}, arXiv:2601.09343 (2026), \url{https://arxiv.org/abs/2601.09343}.
\end{thebibliography}
\end{document}
